\section{Exact full Jones computation with binary planar tensors}
\label{sec:tensor}

The extra logarithm in the faithful partition backend comes from remembering
an equality relation among frontier objects. A binary tensor model can instead
remember one orientation bit per frontier edge. The difficulty is the loop
factor: summing the two orientations of a smoothed circle must produce the
Kauffman factor, independently of its shape or nesting. A combinatorial
rotation cochain supplies this information without drawing coordinates or
floating-point angles.

\subsection{A checked rotation cochain}

Assume first that $n>0$. The validated one-component classical diagram gives
a connected four-valent graph cellularly embedded in the sphere. Multiple
edges, loops, and repeated vertices in a face boundary are allowed. Its
$4n$ darts are numbered by crossing and counterclockwise port. Let $\alpha$
reverse an edge and let $\rho$ advance one port at a crossing. The face walks
are the cycles of $\rho\alpha$, in the convention used by the implementation.
There are $F=n+2$ faces and $2n$ edges. A distinguished face $o$ is regarded as
the outer face.

Assign an integer bend $\beta(d)$ to each dart, with
$\beta(\alpha d)=-\beta(d)$. The local corner of a positive face-boundary
walk contributes one quarter turn. We require
\begin{equation}
 \sum_{d\in f}\beta(d)=b_f,\qquad
 b_f=\begin{cases}4-|f|,&f\ne o,\\-4-|o|,&f=o.\end{cases}
 \label{eq:face-bends}
\end{equation}
Thus bounded faces have total turn $4$ and the outer face has total turn
$-4$. The word ``positive'' here denotes the selected permutation convention;
no additional claim that a face lies on the left of its traversal is needed.

The demand sum is zero:
\[
 \sum_f b_f=4(F-1)-4-4n=0.
\]
Root a spanning tree of the dual graph at $o$, and send each subtree's total
demand through its parent edge. Put zero bend on non-tree dual edges. This
constructs an integral antisymmetric solution. The implementation checks
every antisymmetry equation and every face sum after routing; it does not
accept the routing procedure itself as the check.

\begin{lemma}[A polynomial bend envelope]
The routed cochain can be chosen to satisfy
\[
 \max_d|\beta(d)|\leq3(n+1),\qquad
 \sum_{\{d,\alpha d\}}|\beta(d)|\leq3(n+1)^2.
\]
\end{lemma}
\begin{proof}
Every face is nonempty. A positive demand is at most $3$, and only the
$n+1$ bounded faces can have positive demand. Since total demand is zero,
the absolute demand of any set of faces is at most total positive demand,
which is at most $3(n+1)$. The dual spanning tree has $F-1=n+1$ edges.
The routed value on each is a subtree demand sum, giving both bounds.
\end{proof}

No metric embedding is required. In particular, long arithmetic coordinates
from a planar straight-line drawing never enter this calculation.

\subsection{Why every smoothed circle has rotation four}

A smoothing at a crossing joins adjacent ports. Write
$\tau(a,b)=+1$ when $b-a\equiv1\pmod4$ and $\tau(a,b)=-1$ when
$b-a\equiv-1\pmod4$. Traversing a smoothing from incoming port $a$ to
outgoing port $b$ contributes $\tau(a,b)$. An edge traversal contributes
the corresponding signed bend. We must prove that the total is always
$+4$ or $-4$ on an oriented smoothed circle.

\begin{lemma}[Rotation of a smoothing component]
For any choice of all $n$ smoothings, every resulting circle has total
combinatorial rotation $+4$ in one orientation and $-4$ in the other.
\label{lem:rotation}
\end{lemma}
\begin{proof}
Work with abstract face disks rather than assuming that each face boundary
is a simple polygon in the projection graph. This retains the correct local
sectors at a loop or a repeated vertex. A chosen smoothing creates a band
between the two uncapped opposite face sectors. On the boundary of the
region receiving the band, two old $+1$ corners become two $-1$ turns. Its
total boundary rotation decreases by $4$, exactly as its Euler characteristic
decreases by $1$ when a band is attached.

Initially each bounded face disk has rotation $4$ by
\eqref{eq:face-bends}. Consequently every bounded complementary region $R$
of the smoothed circles has total boundary rotation $4\chi(R)$. This
argument permits bands within one region and between distinct regions; it
depends only on the Euler change under a band attachment.

An innermost smoothed circle bounds a disk and therefore has rotation $4$
in its positive boundary direction. Now proceed outward through the nesting
of circles. A bounded complementary region with $h$ holes has
$\chi(R)=1-h$. Each already treated inner boundary contributes $-4$ with
the region's boundary direction. Therefore its outer boundary has rotation
$4(1-h)+4h=4$. This proves the assertion for all circles. Reversing direction
negates every corner turn and bend, giving $-4$.
\end{proof}

This proof addresses precisely the cases that a naive appeal to a polygonal
turning-number picture can overlook: repeated face walks, self-loops, and
nested components. It also explains why assigning a formal loop weight at
an arbitrary frontier connectivity is unnecessary in the binary model.

\subsection{Local tensors and normalization}

Use the Kauffman bracket convention in which a state with $c$ circles has
loop factor $\delta^c$, with
\[
 \delta=-A^2-A^{-2}.
\]
The normalized Jones polynomial divides by one loop factor. Make the formal
substitution
\[
 A=i z^2,\qquad \delta=z^4+z^{-4}.
\]
Factor out one $i$ per crossing. Smoothing $s=0$ then has weight $z^2$,
and smoothing $s=1$ has weight $-z^{-2}$. All remaining arithmetic is in
$\Z[z,z^{-1}]$.

At a crossing let $\varepsilon_j\in\{+1,-1\}$ denote the orientation of
port $j$, positive when pointing out of the crossing. Edge contraction
identifies opposite port signs. A smoothing joins opposite signs, so a
nonzero crossing tensor has exactly two positive signs and two negative
signs: at most six assignments. The pairings are
\[
 \mathcal P_0=\{(0,1),(2,3)\},\qquad
 \mathcal P_1=\{(0,3),(1,2)\}.
\]
For a compatible pair $(a,b)$ the directed corner contribution is
$-\varepsilon_a\tau(a,b)$, because $\varepsilon_a=-1$ means that the
traversal enters at $a$. Absorb an edge bend at the smaller-numbered dart
only. If $E_v(\varepsilon)$ is the sum of the absorbed signed bends at
crossing $v$, the local tensor entry is
\begin{equation}
 T_v(\varepsilon)=
 \sum_{\substack{s\in\{0,1\}\\
           \varepsilon_a=-\varepsilon_b\ ((a,b)\in\mathcal P_s)}}
 (-1)^s z^{\,2-4s+E_v(\varepsilon)
                   -\sum_{(a,b)\in\mathcal P_s}\varepsilon_a\tau(a,b)}.
 \label{eq:local-tensor}
\end{equation}
Each entry has at most two signed monomials. A loop edge at one crossing
enforces its two endpoint signs locally and never creates a frontier edge.

For clarity, the entries before multiplication by $z^{E_v}$ are as follows.
The asymmetry comes from the fixed turning convention; the closed invariant
is independent of it.
\begin{center}
\begin{tabular}{@{}cccccc@{}}
\toprule
$(+,+,-,-)$ & $(+,-,+,-)$ & $(-,+,+,-)$ & $(+,-,-,+)$ & $(-,+,-,+)$ & $(-,-,+,+)$\\
\midrule
$-z^{-2}$ & $0$ & $z^2$ & $z^2$ & $z^4-z^{-4}$ & $-z^{-2}$\\
\bottomrule
\end{tabular}
\end{center}
The upper bound six is deliberately conservative and remains valid when
local cancellations or loop constraints remove assignments.

Let $F_D(z)$ be the contraction of these tensors. Expanding each local
smoothing choice and then each compatible orientation orients every
smoothed circle independently. Lemma~\ref{lem:rotation} makes its two
orientation weights sum to $z^4+z^{-4}$. Therefore
\begin{equation}
 \langle D\rangle\big|_{A=i z^2}=i^nF_D(z).
 \label{eq:bracket-tensor}
\end{equation}
If $w(D)$ is the writhe, the usual oriented normalization gives
\begin{equation}
 V_D(z^{-8})=
 (-1)^{(n-w(D))/2}z^{-6w(D)}
 \frac{F_D(z)}{z^4+z^{-4}}.
 \label{eq:jones-normalization}
\end{equation}
The exponent $(n-w(D))/2$ is integral. The Laurent implementation performs
an exact division with a checked zero remainder, applies the sign and shift,
and checks that every remaining $z$-exponent is divisible by eight. It
publishes the coefficient dictionary in variable $t=z^{-8}$. The empty
crossing diagram is treated explicitly as the one-component unknot.

\begin{proposition}[Conservative output bounds]
For an $n$-crossing knot diagram,
\[
 V_D(t)=\sum_e c_e t^e,\qquad |e|\leq2n,
 \qquad\sum_e|c_e|\leq4^n.
 \label{prop:jones-bounds}
\]
\end{proposition}
\begin{proof}
In the state sum after division by one loop factor, a state contributes a
monomial times $\delta^{c-1}$. There are $2^n$ states and at most $n+1$
circles in a state of a connected $n$-crossing projection. Its coefficient
mass is at most $2^{c-1}\leq2^n$, so the sum has mass at most $4^n$.
The smoothing monomial has $A$-degree between $-n$ and $n$; the loop factor
adds at most $2n$ in absolute degree, and the writhe correction adds at most
$3n$. Thus the normalized $A$-degree is at most $6n$ in absolute value.
The standard knot normalization has integral $t=A^{-4}$ exponents, which
gives the more than sufficient bound $|e|\leq2n$.
\end{proof}

The code also verifies $V_D(1)=1$ and these bounds. These are useful error
checks, not a substitute for the state-sum proof or an independent polynomial
oracle. A wrong polynomial can satisfy all three.

\subsection{Frontier contraction and the square-root bound}

Fix an ordering of the crossings. After a prefix has been processed, let
$\partial S$ be the edges with exactly one processed endpoint. A state is
one bit for each edge of $\partial S$, with the bit referenced to its
canonical dart. The coefficient for that bit assignment is the sum of all
compatible tensor products in the processed part. The update joins the
current crossing to every compatible assignment, forgets closed edges,
and introduces the new frontier edges. It sums equal output assignments.
There is no remembered partition and no assumption that the processed part
is a disk.

Let $w$ be the maximum number of crossing edges in a prefix cut. There are
at most $2^w$ keys and at most six local transitions per incoming key.
Indexing the bits costs only polynomial overhead. At any intermediate stage
the formal exponent envelope is
\begin{equation}
 |e|\leq B_n:=4n+3(n+1)^2.
 \label{eq:tensor-exponents}
\end{equation}
Indeed a crossing contributes a smoothing exponent of absolute value at most
two and two corner turns of total absolute value at most two; the absorbed
bends contribute at most the cochain's total absolute bend. Expanding at
most six assignments and two monomials per crossing bounds the total
coefficient mass of any partial sum by $12^n$. A coefficient therefore has
$O(n)$ bits, and there are $O(n^2)$ possible Laurent exponents per key.

\begin{theorem}[Binary frontier full-Jones bound]
Given a validated $n$-crossing knot diagram and an order of crossing-edge
frontier width $w$, the full Jones polynomial can be computed exactly in
$\poly(n)2^w$ bit time and space. With the repository's checked planar
separator order this becomes
\[
 \poly(n)2^{O(\sqrt n)}.
\]
\label{thm:tensor}
\end{theorem}
\begin{proof}
The invariant of the update is the partial contraction just described;
induction on the prefix proves it. At the end the frontier is empty, so
\eqref{eq:bracket-tensor} and \eqref{eq:jones-normalization} give the result.
The state count and polynomial arithmetic bounds prove the first claim.

For completeness, a bounded-degree planar graph has recursively balanced
separators of size $O(\sqrt m)$ on an $m$-vertex subproblem
\cite{lipton1979}. List each child subproblem before its separator, as in
the maintained constructor. At a
prefix inside a child, a crossing edge leaving the prefix is incident to
one of the ancestor separators or the current leaf. Since degree is four,
the frontier is bounded by four times the sum of the separator sizes on
that recursion path. Balance gives a geometric series
$\sqrt n+\sqrt{\theta n}+\sqrt{\theta^2n}+\cdots=O(\sqrt n)$ for
some fixed $\theta<1$. The maintained separator constructor supplies an
explicit partition hierarchy; its verifier checks separation, balance,
size, the flattened order, and the resulting width. The tensor backend
constructs and checks this witness before relying on its guarantee.
\end{proof}

The implementation can compare this order with an ordinary scan order and
use the narrower one; the certified witness still bounds the chosen width.
An explicit \code{certified=False} call is a useful ordering ablation, but
it has no universal square-root width promise. Supplying a bad order does
not invalidate the invariant, only its performance.

\subsection{Exact specialization and balanced-digit recovery}

Python sparse Laurent dictionaries carry appreciable overhead. We can
evaluate the tensors at one sufficiently large exact integer and reconstruct
the entire final polynomial, retaining the same frontier state count. This
is a coefficient-bounded Kronecker substitution, not a numerical sample used
as probabilistic evidence.

Choose
\[
 k=\max\!\left(1,\left\lceil\frac{2n+2}{8}\right\rceil\right),
 \qquad z=2^k,\qquad x=z^8\geq4\cdot4^n.
\]
For a first implementation, subtract each local tensor's minimum exponent
$g_v$ and contract its integral value. If $g=\sum_v g_v$, the final integer
is $Z=z^{-g}F_D(z)$. Equation~\eqref{eq:jones-normalization} yields
\begin{equation}
 P(x):=x^{2n}V_D(x^{-1})=
 \frac{(-1)^{(n-w)/2}z^{16n+4-6w+g}Z}{z^8+1}.
 \label{eq:integer-jones}
\end{equation}
A negative power of $z$ is moved into the denominator. Exact integer division
must have zero remainder. The polynomial $P$ has degree at most $4n$, and
each coefficient has absolute value at most $4^n<x/2$.

\begin{lemma}[Balanced recovery]
An integer of the form $P(x)=\sum_{j=0}^{4n}a_jx^j$, with
$|a_j|<x/2$, determines all the $a_j$ uniquely. Iteratively choose the
residue in $[-x/2,x/2)$ and divide the remaining difference by $x$.
\end{lemma}
\begin{proof}
The constant coefficient is the unique allowed representative of the
integer modulo $x$. Subtracting it and dividing by $x$ leaves the same
problem with one fewer coefficient. Induction proves uniqueness and the
algorithm. The recovered coefficient at index $j$ is the coefficient of
$t^{2n-j}$ in $V_D(t)$.
\end{proof}

The intermediate specialized tensor values need not encode their formal
Laurent polynomials injectively. Distinct partial polynomials can specialize
to the same value, and a nonzero partial polynomial can specialize to zero.
Evaluation is a ring homomorphism, so every subsequent addition and
multiplication still gives the correct final value. Faithfulness is needed
only for final recovery, where the proved degree and coefficient bounds
apply. This distinction permits much smaller intermediate keys and avoids
an unjustified assumption about partial coefficient growth.

The implementation includes an exact identity shortcut: compare the
contracted value to the value implied by $V_D=1$ before normalization and
digit extraction when its monomial exponent permits that comparison. The
same faithful encoding proves this equality is a whole-polynomial identity
test. It does not make the invariant a complete unknot certificate.

\subsection{Valuation normalization and gauge covariance}

The global-shift version can spend most of its time carrying powers of $z$
that merely encode monomial offsets. Store a nonzero specialized value in
the canonical form
\[
 z^v m,\qquad v\in\Z,\quad m\in\Z\setminus\{0\},\quad z\nmid m.
\]
Zero has a separate representation. Multiplying by a signed local monomial
adds to $v$ and changes the sign of $m$. To add $z^{v_1}m_1$ and
$z^{v_2}m_2$, put $r=\min(v_1,v_2)$ and form
\[
 a=m_1z^{v_1-r}+m_2z^{v_2-r}.
\]
If $a\ne0$, remove the largest power $z^j$ dividing $a$ and return
$(r+j,a/z^j)$. With $z=2^k$, this uses the number of trailing zero bits of
$|a|$, divided by $k$. It is exact even when carries create additional
factors of $z$. These are valuations of the evaluated rational number,
not formal lowest exponents of an unexpanded polynomial.

At the final state let $F_D(z)=z^v m$. Equation~\eqref{eq:integer-jones}
applies with $g=v$ and $Z=m$, and the same balanced decoding recovers the
full polynomial. From \eqref{eq:tensor-exponents} and the $12^n$ mass bound,
\[
 -B_n\leq v\leq B_n+\log_z(12^n),\qquad
 \bits(m)\leq2kB_n+n\log_2 12+O(1)=O(n^3).
\]
Thus arbitrary-precision integer work remains polynomial per transition.
The bound is deliberately conservative; the experiments record actual
mantissa sizes.

\begin{proposition}[Gauge covariance of partial contractions]
Fix the same outer face. If $\beta$ and $\beta'$ both satisfy
\eqref{eq:face-bends}, their difference is an integral vertex coboundary
on the projection graph. For each fixed frontier assignment, the two partial
contraction coefficients differ by a monomial $z^q$ depending only on that
assignment and the processed crossing set. Closed contractions agree.
Consequently the normalized mantissas and their bit lengths agree.
\label{prop:gauge}
\end{proposition}
\begin{proof}
The difference has zero period on every facial cycle. Since the embedding
is spherical, facial cycles generate the integral cycle space, so integrating
along a primal spanning tree yields an integral potential $h$ with
$\beta'-\beta=\delta h$. In a tensor product, the resulting exponent is a
sum of signed edge potential differences. At each internal crossing, the
two-in/two-out condition cancels the coefficient of $h$. Only boundary
terms remain; the convention for absorbing a bend at one endpoint fixes
which boundary term is used. All internal assignments with the same
frontier signs therefore receive the same monomial multiplier. It factors
through every partial sum. At a closed frontier the multiplier is one.
Canonical valuation normalization removes precisely such a multiplier
from the mantissa.
\end{proof}

This also gives a useful kernel-level robustness statement. An arbitrary
valid same-outer cochain with $L$-bit bends can have enormous numerical
values. Its boundary gauge offsets need only $O(L+\log(n+2))$ bits; the
mantissas and the differences used to align terms are those of the bounded
canonical cochain. A valuation implementation therefore has
$\poly(n,L)2^w$ cost on such data. The public constructor currently creates
the bounded canonical cochain; accepting externally supplied gauges is
not a new public interface in this patch.

Gauge covariance alone does \emph{not} prove that every intermediate formal
Laurent span is $O(n)$. It removes a common monomial offset, not the span
between genuinely different exponents. We retain the proved $O(n^2)$
envelope for the complexity theorem and leave sharper span bounds as a
separate research problem.

\subsection{Recognition semantics and limits}

A completed polynomial different from one proves knottedness. A completed
identity polynomial is retained as evidence and the recognizer continues to
its other stages. An exhausted state, transition, or global time allowance
publishes no partial polynomial as an invariant. This distinction is checked
both through the standalone Jones command and through the public recognizer.

The patch introduces \code{jones\_backend="tensor"} and the CLI choice
\code{--backend tensor}. The optional public dispatch uses valuation-normalized
integer arithmetic. The low-level \code{tensor\_jones} function defaults to
the Laurent implementation as a readable reference and supports the two
integer modes explicitly. The existing default Jones backend is unchanged.
The square-root theorem concerns uncapped complete invariant computation;
capped queries are partial decision attempts and may be inconclusive.
